\subsection{Periodic-boundary ground space for a block-diagonal tensor}
\label{subsec:block_diagonal_chain_equality}

The periodic ground-space theorem states that, for sufficiently long periodic
systems, the periodic-boundary ground space of a parent Hamiltonian is generated
by a basis of normal tensors of the initial tensor
$A$~\cite[Theorem~IV.16]{Cirac2021Matrix}. The lemmas in this subsection are
conditional reductions toward that statement. Their
hypotheses explicitly include a boundary inclusion, a block-diagonal boundary
representation, or periodic-boundary membership for the single-block states; in
the source theorem these conditions are obtained from the boundary-condition
comparison in the periodic ring. The source proof
\cite{PerezGarcia2007Matrix} writes this comparison with $C^j_{i_1}$,
$D^j_{i_{m+1}}$, and the normalized matrix $E^j$; the symbols $\beta$ and
$\rho$ below are the words obtained after opening the periodic boundary at the
chosen cut. The source theorem writes the number of blocks as $b$ and assumes
$b\ge2$; here the same number is denoted by $g$. In these coordinates the
source comparison input becomes
$A^j_\beta C^j_{i,\rho}=(\mu_j^N X_j A^j_\beta)A^j_\rho$,
which is the form obtained after opening the boundary from
$A^j_{i_{m+1}}C^j_{i_1}=D^j_{i_{m+1}}A^j_{i_1}$ with
$D^j_\beta=(\mu_j^N X_j)A^j_\beta$.
Once this compatibility is known, the normalized $E^j$ calculation is the
matrix identity of
Lemma~\ref{lem:block_boundary_cde_identities_from_trace_decompositions};
the local boundary form used after opening the periodic boundary is
Lemma~\ref{lem:complementary_word_boundary_identities_from_pgvwc}.

\begin{lemma}[The block-diagonal periodic-boundary space lies between two block sums]
    \label{lem:bnt_block_diagonal_periodic_chain_inclusions}
    \lean{MPSTensor.chainGroundSpace_toTensorFromBlocks_two_inclusions_and_iSupIndep_of_bnt_unital_c1}
    \leanok
    \uses{lem:isup_chain_ground_space_block_le_assembled,
        lem:bnt_block_diagonal_periodic_constraints_internal_sum}
    With the hypotheses and notation of
    Lemma~\ref{lem:bnt_block_diagonal_periodic_constraints_propagated_sum}, one has
    \begin{align}
        \bigvee_{j=1}^g\mathcal G_{N,L}(A_j)
        &\subseteq
        \mathcal G_{N,L}\!\left(\bigoplus_{j=1}^g\mu_jA_j\right) \subseteq
        \bigvee_{j=1}^gG_N(A_j),
        \notag
    \end{align}
    and the sum defining the right-hand side is internal.
\end{lemma}

\begin{proof}
    \leanok
    The first inclusion is
    Lemma~\ref{lem:isup_chain_ground_space_block_le_assembled}. The second
    inclusion and internal directness are
    Lemma~\ref{lem:bnt_block_diagonal_periodic_constraints_internal_sum}.
    % The remaining arXiv:quant-ph/0608197 periodic-boundary comparison with
    % block-diagonal boundary
    % conditions replaces
    % $\bigvee_jG_N(A_j)$ by $\sum_j\mathcal G_{N,L}(A_j)$.
\end{proof}

\begin{lemma}[Boundary inclusion gives the block-diagonal periodic-boundary equality]
    \label{lem:block_diagonal_boundary_closing_equality}
    \lean{MPSTensor.chainGroundSpace_toTensorFromBlocks_eq_iSup_chainGroundSpace_of_boundary_closing}
    \leanok
    \uses{lem:isup_chain_ground_space_block_le_assembled}
    Let $B=\bigoplus_{j=1}^g\mu_jA_j$, where $\mu_j\ne0$ for
    $1\le j\le g$. If the periodic-boundary comparison with block-diagonal
    boundary conditions gives
    \begin{align}
        \mathcal G_{N,L}(B)
        &\subseteq
        \bigvee_{j=1}^g\mathcal G_{N,L}(A_j),
        \notag
    \end{align}
    then
    $\mathcal G_{N,L}(B) = \bigvee_{j=1}^g\mathcal G_{N,L}(A_j)$.
\end{lemma}

\begin{proof}
    \leanok
    The reverse inclusion is
    Lemma~\ref{lem:isup_chain_ground_space_block_le_assembled}; combine the two
    inclusions.
\end{proof}

\begin{lemma}[Componentwise periodic-boundary decomposition gives the reverse inclusion]
    \label{lem:block_diagonal_boundary_decomposition_inclusion}
    \lean{MPSTensor.chainGroundSpace_toTensorFromBlocks_le_iSup_of_boundary_decomposition}
    \leanok
    Let $B=\bigoplus_{j=1}^g\mu_jA_j$. Suppose that every vector
    $\psi\in\mathcal G_{N,L}(B)$ can be written as
    \begin{align}
        \psi &= \sum_{j=1}^g\psi_j,
        \notag\\
        \psi_j &\in \mathcal G_{N,L}(A_j).
        \notag
    \end{align}
    Then
    \begin{align}
        \mathcal G_{N,L}(B)
        &\subseteq
        \bigvee_{j=1}^g\mathcal G_{N,L}(A_j).
        \notag
    \end{align}
\end{lemma}

\begin{proof}
    \leanok
    For $\psi\in\mathcal G_{N,L}(B)$, choose the displayed decomposition. Each
    summand lies in the corresponding subspace
    $\mathcal G_{N,L}(A_j)$, hence their finite sum lies in
    $\bigvee_j\mathcal G_{N,L}(A_j)$.
\end{proof}

\begin{lemma}[Block-diagonal boundary vectors in the periodic-boundary block sum]
    \label{lem:block_diagonal_boundary_vector_mem_chain_sum}
    \lean{MPSTensor.groundSpaceMap_toTensorFromBlocks_blockDiagonal_mem_iSup_chainGroundSpace}
    \leanok
    \uses{lem:block_diagonal_boundary_condition_sum}
    Let $B=\bigoplus_{j=1}^g\mu_jA_j$. Fix block-diagonal boundary conditions
    $X_j$. Suppose that, for every $j$,
    $\Gamma_N^{A_j}(\mu_j^N X_j)\in\mathcal G_{N,L}(A_j)$. Then
    \begin{align}
        \Gamma_N^B\!\left(\bigoplus_{j=1}^gX_j\right)
        &\in
        \bigvee_{j=1}^g\mathcal G_{N,L}(A_j).
        \notag
    \end{align}
\end{lemma}

\begin{proof}
    \leanok
    The block-diagonal boundary formula gives
    \begin{align}
        \Gamma_N^B\!\left(\bigoplus_{j=1}^gX_j\right)
        &=
        \sum_{j=1}^g\Gamma_N^{A_j}(\mu_j^N X_j).
        \notag
    \end{align}
    Each summand belongs to the corresponding block periodic-boundary space by
    assumption, so the finite sum lies in the linear span of these spaces.
\end{proof}

\begin{lemma}[Block-diagonal boundary representation gives the reverse inclusion]
    \label{lem:block_diagonal_boundary_representation_inclusion}
    \lean{MPSTensor.chainGroundSpace_toTensorFromBlocks_le_iSup_of_blockDiagonal_boundary_groundSpaceMap}
    \leanok
    \uses{lem:block_diagonal_boundary_vector_mem_chain_sum}
    Let $B=\bigoplus_{j=1}^g\mu_jA_j$. Suppose every
    $\psi\in\mathcal G_{N,L}(B)$ admits block-diagonal boundary conditions
    $X_j$ such that
    $\psi = \Gamma_N^B\!\left(\bigoplus_{j=1}^gX_j\right)$,
    and, for every $j$,
    $\Gamma_N^{A_j}(\mu_j^N X_j)\in\mathcal G_{N,L}(A_j)$. Then
    \begin{align}
        \mathcal G_{N,L}(B)
        &\subseteq
        \bigvee_{j=1}^g\mathcal G_{N,L}(A_j).
        \notag
    \end{align}
\end{lemma}

\begin{proof}
    \leanok
    Apply Lemma~\ref{lem:block_diagonal_boundary_vector_mem_chain_sum} to the
    block-diagonal boundary representation of each
    $\psi\in\mathcal G_{N,L}(B)$.
\end{proof}

\begin{lemma}[Block-diagonal boundary representation reduces to periodic-boundary equality]
    \label{lem:bnt_block_diagonal_boundary_representation_equality}
    \lean{MPSTensor.chainGroundSpace_toTensorFromBlocks_eq_iSup_and_iSupIndep_of_bnt_c1_blockBoundary}
    \leanok
    \uses{lem:block_diagonal_boundary_representation_inclusion,
        lem:block_diagonal_boundary_closing_equality,
        lem:bnt_block_diagonal_periodic_constraints_internal_sum}
    With the hypotheses and notation of
    Lemma~\ref{lem:bnt_block_diagonal_periodic_constraints_propagated_sum}, suppose
    every vector $\psi\in\mathcal G_{N,L}(B)$ has block-diagonal boundary
    conditions $X_j$ such that
    $\psi = \Gamma_N^B\!\left(\bigoplus_{j=1}^gX_j\right)$,
    and, for every $j$,
    $\Gamma_N^{A_j}(\mu_j^N X_j)\in\mathcal G_{N,L}(A_j)$. Then
    $\mathcal G_{N,L}(B) = \bigvee_{j=1}^g\mathcal G_{N,L}(A_j)$,
    and the sum defining $\bigvee_jG_N(A_j)$ is internal.
\end{lemma}

\begin{proof}
    \leanok
    The block-diagonal boundary representation gives
    \begin{align}
        \mathcal G_{N,L}(B)
        &\subseteq
        \bigvee_{j=1}^g\mathcal G_{N,L}(A_j)
        \notag
    \end{align}
    by Lemma~\ref{lem:block_diagonal_boundary_representation_inclusion}. The
    equality follows from Lemma~\ref{lem:block_diagonal_boundary_closing_equality};
    the internality statement is
    Lemma~\ref{lem:bnt_block_diagonal_periodic_constraints_internal_sum}.
\end{proof}

\begin{lemma}[Complementary-word matrix identities reduce to periodic-boundary equality]
    \label{lem:bnt_block_diagonal_complementary_identities_equality}
    \lean{MPSTensor.chainGroundSpace_toTensorFromBlocks_eq_iSup_and_iSupIndep_of_complementary_identities}
    \leanok
    \uses{lem:block_diagonal_boundary_periodic_components_from_complementary_identities,
        lem:bnt_block_diagonal_boundary_representation_equality}
    With the hypotheses and notation of
    Lemma~\ref{lem:bnt_block_diagonal_boundary_representation_equality}, assume
    in addition that $L+L_0\le N$. Assume further that, for every
    $\psi\in\mathcal G_{N,L}\!\left(\bigoplus_j\mu_jA_j\right)$
    and every block-diagonal boundary representation
    \begin{align}
        \psi
        &=
        \Gamma_N^{\oplus_j\mu_jA_j}\!\left(\bigoplus_jX_j\right),
        \notag
    \end{align}
    the complementary-word matrix identities hold: for every block $j$,
    boundary-crossing interval $i$, word $\beta$ before the cut, and outside
    word $\rho$ of length $N-L$, there is a matrix $E_{j,i,\rho}$ such that
    $\mu_j^N X_jA^j_\beta A^j_\rho = A^j_\beta E_{j,i,\rho}$.
    Then
    \begin{align}
        \mathcal G_{N,L}\!\left(\bigoplus_j\mu_jA_j\right)
        &=
        \bigvee_j\mathcal G_{N,L}(A_j),
        \notag
    \end{align}
    and the $N$-site ground spaces $G_N(A_j)=\mathcal G_{N,N}(A_j)$
    have internal sum.
\end{lemma}

\begin{proof}
    \leanok
    \uses{lem:block_diagonal_boundary_periodic_components_from_complementary_identities,
        lem:bnt_block_diagonal_boundary_representation_equality}
    For each $\psi$,
    Lemma~\ref{lem:block_diagonal_boundary_periodic_components_from_complementary_identities}
    gives matrices $X_j$ with
    \begin{align}
        \psi
        &=
        \Gamma_N^{\oplus_j\mu_jA_j}\!\left(\bigoplus_jX_j\right)
        \notag
    \end{align}
    and, for every $j$, the membership
    $\Gamma_N^{A_j}(\mu_j^N X_j)\in\mathcal G_{N,L}(A_j)$ holds.
    Lemma~\ref{lem:bnt_block_diagonal_boundary_representation_equality}
    then gives
    \begin{align}
        \mathcal G_{N,L}\!\left(\bigoplus_j\mu_jA_j\right)
        &=
        \bigvee_j\mathcal G_{N,L}(A_j),
        \notag
    \end{align}
    and the internality of the sum of the $G_N(A_j)$.
\end{proof}

\begin{lemma}[$C^j,D^j$ boundary-condition comparisons reduce to periodic-boundary equality]
    \label{lem:bnt_block_diagonal_pgvwc_comparison_equality}
    \lean{MPSTensor.chainGroundSpace_toTensorFromBlocks_eq_iSup_and_iSupIndep_of_pgvwc_comparison}
    \leanok
    \uses{lem:block_diagonal_boundary_periodic_components_from_pgvwc_comparison,
        lem:bnt_block_diagonal_boundary_representation_equality}
    With the hypotheses and notation of
    Lemma~\ref{lem:bnt_block_diagonal_boundary_representation_equality}, assume
    in addition that $L+L_0\le N$. Assume further that, for every
    $\psi\in\mathcal G_{N,L}\!\left(\bigoplus_j\mu_jA_j\right)$
    and every block-diagonal boundary representation
    \begin{align}
        \psi
        &=
        \Gamma_N^{\oplus_j\mu_jA_j}\!\left(\bigoplus_jX_j\right),
        \notag
    \end{align}
    there are matrices $C^j_{i,\rho}$ such that, for every
    boundary-crossing interval $i$, every outside word $\rho$ of length $N-L$,
    and every word $\beta$ before the cut,
    $A^j_\beta C^j_{i,\rho} = ((\mu_j^N X_j)A^j_\beta)A^j_\rho$.
    This is the $C^j,D^j$ compatibility input; the subsequent normalized
    $E^j$ identities follow from
    Lemma~\ref{lem:complementary_word_boundary_identities_from_pgvwc}.
    Then
    \begin{align}
        \mathcal G_{N,L}\!\left(\bigoplus_j\mu_jA_j\right)
        &=
        \bigvee_j\mathcal G_{N,L}(A_j),
        \notag
    \end{align}
    and the $N$-site ground spaces $G_N(A_j)=\mathcal G_{N,N}(A_j)$
    have internal sum.
\end{lemma}

\begin{proof}
    \leanok
    \uses{lem:block_diagonal_boundary_periodic_components_from_pgvwc_comparison,
        lem:bnt_block_diagonal_boundary_representation_equality}
    For each $\psi$,
    Lemma~\ref{lem:block_diagonal_boundary_periodic_components_from_pgvwc_comparison}
    gives matrices $X_j$ with
    \begin{align}
        \psi
        &=
        \Gamma_N^{\oplus_j\mu_jA_j}\!\left(\bigoplus_jX_j\right)
        \notag
    \end{align}
    and, for every $j$,
    $\Gamma_N^{A_j}(\mu_j^N X_j)\in\mathcal G_{N,L}(A_j)$.
    Lemma~\ref{lem:bnt_block_diagonal_boundary_representation_equality}
    then gives
    \begin{align}
        \mathcal G_{N,L}\!\left(\bigoplus_j\mu_jA_j\right)
        &=
        \bigvee_j\mathcal G_{N,L}(A_j),
        \notag
    \end{align}
    and the internality of the sum of the $G_N(A_j)$.
\end{proof}

\begin{lemma}[$C^j,D^j$ comparisons give the periodic block ground-space equality]
    \label{lem:bnt_block_diagonal_pgvwc_comparison_ground_space}
    \lean{MPSTensor.chainGroundSpace_toTensorFromBlocks_eq_iSup_of_pgvwc_comparison}
    \leanok
    \uses{lem:bnt_block_diagonal_pgvwc_comparison_equality}
    With the hypotheses and notation of
    Lemma~\ref{lem:bnt_block_diagonal_pgvwc_comparison_equality}, one has
    \begin{align}
        \mathcal G_{N,L}\!\left(\bigoplus_j\mu_jA_j\right)
        &=
        \bigvee_j\mathcal G_{N,L}(A_j).
        \notag
    \end{align}
    This records only the ground-space equality conclusion obtained after the
    $C^j,D^j$ boundary-condition comparison input in
    \cite[Theorem~12, proof lines~1446--1456]{PerezGarcia2007Matrix}; the
    local normalized $E^j$ calculation is
    Lemma~\ref{lem:complementary_word_boundary_identities_from_pgvwc}.
    The independence of the $N$-site spaces is a separate conclusion of
    Lemma~\ref{lem:bnt_block_diagonal_pgvwc_comparison_equality}.
    The statement is still in the length-$L_0$ injectivity range displayed in
    that lemma, not the shorter range in
    \cite[Theorem~12]{PerezGarcia2007Matrix}.
\end{lemma}

\begin{proof}
    \leanok
    The claim is the equality part of
    Lemma~\ref{lem:bnt_block_diagonal_pgvwc_comparison_equality}.
\end{proof}

\begin{lemma}[$C^j,D^j$ comparisons at boundary-crossing intervals reduce to periodic-boundary equality]
    \label{lem:bnt_block_diagonal_crossing_pgvwc_comparison_equality}
    \lean{MPSTensor.chainGroundSpace_toTensorFromBlocks_eq_iSup_and_iSupIndep_of_crossing_pgvwc_comparison}
    \leanok
    \uses{lem:bnt_block_diagonal_boundary_representation_equality,
        thm:block_diagonal_boundary_crossing_pgvwc_comparison_chain_ground_space,
        lem:block_diagonal_boundary_component_pgvwc_comparison_injective}
    With the hypotheses and notation of
    Lemma~\ref{lem:bnt_block_diagonal_boundary_representation_equality}, assume
    in addition that $L+L_0\le N$. Assume also that, for every
    boundary-crossing interval beginning at $i$, the simultaneous products
    \begin{align}
        w &\longmapsto (A^1_w,\ldots,A^g_w),
        \notag\\
        |w| &= N-i,
        \notag
    \end{align}
    span the product algebra $\prod_j\MN{D_j}$. Then
    \begin{align}
        \mathcal G_{N,L}\!\left(\bigoplus_j\mu_jA_j\right)
        &=
        \bigvee_j\mathcal G_{N,L}(A_j),
        \notag
    \end{align}
    and the $N$-site ground spaces $G_N(A_j)=\mathcal G_{N,N}(A_j)$
    have internal sum.
    The present statement is the span-dependent consequence, after opening the
    boundary, of the $C^j,D^j$ compatibility input in
    \cite[Theorem~12, proof lines~1436--1451]{PerezGarcia2007Matrix} in the
    length-$L_0$ injectivity range
    $L \ge (L_0+1)+3(g-1)(L_0+1)+1$.
    The source theorem has the shorter bound
    $L \ge 3(g-1)(L_0+1)+1$.
\end{lemma}

\begin{proof}
    \leanok
    \uses{lem:bnt_block_diagonal_boundary_representation_equality,
        thm:block_diagonal_boundary_crossing_pgvwc_comparison_chain_ground_space,
        lem:block_diagonal_boundary_component_pgvwc_comparison_injective}
    Lemma~\ref{lem:bnt_block_diagonal_boundary_representation_equality} reduces
    the claim to producing, for each $\psi$, block-diagonal boundary conditions
    $X_j$ whose single-block components lie in
    $\mathcal G_{N,L}(A_j)$. The crossing comparison theorem supplies this:
    it first uses the boundary-comparison-free BNT representation to obtain
    $X_j$, and then uses the stated tail-word spans and
    Theorem~\ref{thm:block_diagonal_boundary_crossing_pgvwc_comparison_chain_ground_space}
    to produce matrices $C^j_{i,\rho}$ satisfying
    $A^j_\beta C^j_{i,\rho} = ((\mu_j^N X_j)A^j_\beta)A^j_\rho$.
    Lemma~\ref{lem:block_diagonal_boundary_component_pgvwc_comparison_injective}
    then gives
    $\Gamma_N^{A_j}(\mu_j^N X_j)\in\mathcal G_{N,L}(A_j)$
    for every $j$. Applying
    Lemma~\ref{lem:bnt_block_diagonal_boundary_representation_equality} to these
    boundary conditions gives the displayed equality and internality.
\end{proof}

\begin{lemma}[Boundary-crossing comparisons give the periodic block ground-space equality]
    \label{lem:bnt_block_diagonal_crossing_pgvwc_comparison_ground_space}
    \lean{MPSTensor.chainGroundSpace_toTensorFromBlocks_eq_iSup_of_crossing_pgvwc_comparison}
    \leanok
    \uses{lem:bnt_block_diagonal_crossing_pgvwc_comparison_equality}
    With the hypotheses and notation of
    Lemma~\ref{lem:bnt_block_diagonal_crossing_pgvwc_comparison_equality}, one has
    \begin{align}
        \mathcal G_{N,L}\!\left(\bigoplus_j\mu_jA_j\right)
        &=
        \bigvee_j\mathcal G_{N,L}(A_j).
        \notag
    \end{align}
    This records the ground-space equality conclusion of the conditional
    crossing-span version above. It is the corresponding part of
    \cite[Theorem~12]{PerezGarcia2007Matrix}, but only under the visible
    tail-word span hypothesis and in the longer stated range. The independence
    of the $N$-site spaces is a separate conclusion of
    Lemma~\ref{lem:bnt_block_diagonal_crossing_pgvwc_comparison_equality}.
    This intermediate statement remains conditional on the short tail-word
    spans. The source-range conclusion without them is supplied by the global
    change-of-cut theorem
    Theorem~\ref{thm:block_diagonal_boundary_periodic_components_from_global_cut}
    only in the currently formalized doubly normalized specialization
    $\Lambda_j=\Id$. For the unrestricted source normalization, with positive
    definite dual fixed points satisfying
    \begin{align}
        \sum_a A_j^{a\,\dagger}\Lambda_j A_j^a&=\Lambda_j, \notag
    \end{align}
    the corresponding source-range theorem remains unformalized.
    % General normalization is tracked in issues #5751 and #5770. Neither issue
    % is cited here as a completed theorem.
\end{lemma}

\begin{proof}
    \leanok
    The claim is the equality part of
    Lemma~\ref{lem:bnt_block_diagonal_crossing_pgvwc_comparison_equality}.
\end{proof}

\begin{lemma}[Boundary trace comparisons reduce to periodic-boundary equality]
    \label{lem:bnt_block_diagonal_trace_decomposition_equality}
    \lean{MPSTensor.chainGroundSpace_toTensorFromBlocks_eq_iSup_and_iSupIndep_of_trace_decomposition}
    \leanok
    \uses{lem:block_diagonal_boundary_periodic_components_from_trace_decomposition,
        lem:bnt_block_diagonal_boundary_representation_equality}
    With the hypotheses and notation of
    Lemma~\ref{lem:bnt_block_diagonal_boundary_representation_equality}, assume
    in addition that $L+L_0\le N$. Assume also that, for some middle-word length
    $m$, the simultaneous tuples $w\mapsto(A^1_w,\ldots,A^g_w)$, where $w$
    ranges over words of length $m$ in $\Fin{d}$, span the full product algebra
    $\prod_j\MN{D_j}$. Assume further that, for every
    $\psi\in\mathcal G_{N,L}\!\left(\bigoplus_j\mu_jA_j\right)$
    and every block-diagonal boundary representation
    \begin{align}
        \psi
        &=
        \Gamma_N^{\oplus_j\mu_jA_j}\!\left(\bigoplus_jX_j\right),
        \notag
    \end{align}
    there are matrices $C^j_{i,\rho}$ such that, for every
    boundary-crossing interval $i$, every outside word $\rho$ of length $N-L$,
    every word $\beta$ before the cut, and every word $w$ of length $m$,
    \begin{align}
        \sum_j\tr\!\left(A^j_\beta C^j_{i,\rho}A^j_w\right)
        &=
        \sum_j
        \tr\!\left(((\mu_j^N X_j)A^j_\beta)A^j_\rho A^j_w\right).
        \notag
    \end{align}
    Then
    \begin{align}
        \mathcal G_{N,L}\!\left(\bigoplus_j\mu_jA_j\right)
        &=
        \bigvee_j\mathcal G_{N,L}(A_j),
        \notag
    \end{align}
    and the $N$-site ground spaces $G_N(A_j)=\mathcal G_{N,N}(A_j)$
    have internal sum.
\end{lemma}

\begin{proof}
    \leanok
    \uses{lem:block_diagonal_boundary_periodic_components_from_trace_decomposition,
        lem:bnt_block_diagonal_boundary_representation_equality}
    For each $\psi$,
    Lemma~\ref{lem:block_diagonal_boundary_periodic_components_from_trace_decomposition}
    gives matrices $X_j$ with
    \begin{align}
        \psi
        &=
        \Gamma_N^{\oplus_j\mu_jA_j}\!\left(\bigoplus_jX_j\right)
        \notag
    \end{align}
    and, for every $j$,
    $\Gamma_N^{A_j}(\mu_j^N X_j)\in\mathcal G_{N,L}(A_j)$.
    Lemma~\ref{lem:bnt_block_diagonal_boundary_representation_equality}
    then gives
    \begin{align}
        \mathcal G_{N,L}\!\left(\bigoplus_j\mu_jA_j\right)
        &=
        \bigvee_j\mathcal G_{N,L}(A_j),
        \notag
    \end{align}
    and the internality of the sum of the $G_N(A_j)$.
\end{proof}

\begin{lemma}[Basis-of-normal-tensors product span gives boundary-trace ground-space equality]
    \label{lem:bnt_block_diagonal_trace_decomposition_bnt_span_equality}
    \lean{MPSTensor.chainGroundSpace_toTensorFromBlocks_eq_iSup_and_iSupIndep_of_trace_decomposition_bnt_c1_span}
    \leanok
    \uses{lem:unital_block_separating_product_span,
        lem:bnt_block_diagonal_trace_decomposition_equality}
    With the hypotheses and notation of
    Lemma~\ref{lem:bnt_block_diagonal_boundary_representation_equality}, assume
    in addition that $L+L_0\le N$ and
    $L \ge (L_0+1)+(g-1)((L_0+1)+((L_0+1)+(L_0+1)))+1$.
    Put
    $m = (L_0+1)+(g-1)((L_0+1)+((L_0+1)+(L_0+1)))$.
    Assume that, for every
    $\psi\in\mathcal G_{N,L}\!\left(\bigoplus_j\mu_jA_j\right)$
    and every block-diagonal boundary representation
    \begin{align}
        \psi
        &=
        \Gamma_N^{\oplus_j\mu_jA_j}\!\left(\bigoplus_jX_j\right),
        \notag
    \end{align}
    there are matrices $C^j_{i,\rho}$ such that, for every
    boundary-crossing interval $i$, every outside word $\rho$ of length $N-L$,
    every word $\beta$ before the cut, and every word $w$ of length $m$,
    \begin{align}
        \sum_j\tr\!\left(A^j_\beta C^j_{i,\rho}A^j_w\right)
        &=
        \sum_j
        \tr\!\left(((\mu_j^N X_j)A^j_\beta)A^j_\rho A^j_w\right).
        \notag
    \end{align}
    Then
    \begin{align}
        \mathcal G_{N,L}\!\left(\bigoplus_j\mu_jA_j\right)
        &=
        \bigvee_j\mathcal G_{N,L}(A_j),
        \notag
    \end{align}
    and the $N$-site ground spaces $G_N(A_j)=\mathcal G_{N,N}(A_j)$
    have internal sum.
\end{lemma}

\begin{proof}
    \leanok
    \uses{lem:unital_block_separating_product_span,
        lem:bnt_block_diagonal_trace_decomposition_equality}
    Lemma~\ref{lem:unital_block_separating_product_span} gives
    $\spn\{(A^1_w,\ldots,A^g_w):|w|=m\} = \prod_j\MN{D_j}$.
    Applying
    Lemma~\ref{lem:bnt_block_diagonal_trace_decomposition_equality} with this
    middle-word span gives the equality of periodic-boundary ground spaces and
    the internality of the sum of the $N$-site spaces.
\end{proof}

\begin{lemma}[Boundary decomposition reduces to periodic-boundary equality]
    \label{lem:bnt_block_diagonal_boundary_decomposition_equality}
    \lean{MPSTensor.chainGroundSpace_toTensorFromBlocks_eq_iSup_and_iSupIndep_of_bnt_c1_boundary_decomposition}
    \leanok
    \uses{lem:block_diagonal_boundary_decomposition_inclusion,
        lem:isup_chain_ground_space_block_le_assembled,
        lem:bnt_block_diagonal_periodic_constraints_internal_sum}
    Let $A_1,\ldots,A_g$ be normalized representatives of a basis of normal
    tensors: each representative is irreducible, the family is left-canonical,
    the self-overlaps are normalized, and no two distinct same-dimensional
    representatives are gauge-phase equivalent, and let
    $B=\bigoplus_{j=1}^g\mu_jA_j$, where $\mu_j\ne0$.
    Assume each block is injective at length $L_0$, assume
    $\sum_aA^j_aA^{j\dagger}_a=\Id$ for every $j$, and suppose
    \begin{align}
        L &\ge (L_0+1)+(g-1)3(L_0+1)+1,
        \notag\\
        N &\ge L.
        \notag
    \end{align}
    If every vector $\psi\in\mathcal G_{N,L}(B)$ can be written as
    \begin{align}
        \psi &= \sum_{j=1}^g\psi_j,
        \notag\\
        \psi_j &\in \mathcal G_{N,L}(A_j),
        \notag
    \end{align}
    then
    $\mathcal G_{N,L}(B) = \bigvee_{j=1}^g\mathcal G_{N,L}(A_j)$,
    and the sum $\bigvee_jG_N(A_j)$ is internal.
\end{lemma}

\begin{proof}
    \leanok
    Lemma~\ref{lem:block_diagonal_boundary_decomposition_inclusion} gives
    \begin{align}
        \mathcal G_{N,L}(B)
        &\subseteq
        \bigvee_{j=1}^g\mathcal G_{N,L}(A_j).
        \notag
    \end{align}
    The blockwise inclusion gives the reverse containment, hence equality.
    Internality follows from the directness conclusion of
    Lemma~\ref{lem:bnt_block_diagonal_periodic_constraints_internal_sum}.
\end{proof}
